\documentclass[12pt]{article}
\usepackage[utf8]{inputenc}
\usepackage[T1]{fontenc}
\usepackage{hyperref}
\title{Magdalena Pe\~nasco Mixtec Language Resources for K-12 Curriculum Development}
\author{Compiled by Claude Sonnet 5.5 (claude-sonnet-5-5) for Andr\'as Kornai}
\date{}
\begin{document}
\maketitle
\section{Introduction}
Magdalena Pe\~nasco Mixtec (ISO 639-3 xtm, Glottocode magd1235) is an Otomanguean language of Oaxaca, Mexico, also called Tlacotepec Mixtec. Sources cite 7,400 speakers (2005 census, with 1,200 monolinguals in 2000) and Joshua Project gives 9,300. Varieties are listed for Santa Mar\'ia Magdalena Pe\~nasco, San Crist\'obal Amoltepec, San Mateo Pe\~nasco and San Agust\'in Tlacotepec.
\section{Primary Educational Resources}
No primers, textbooks or school programmes were found in the search results.
\section{Linguistic and Reference Materials}
\begin{itemize}
\item Wikipedia article on Pe\~nasco Mixtec: \url{https://en.wikipedia.org/wiki/Pe%C3%B1asco_Mixtec}
\item Wiktionary category of xtm terms from Oaxaca: \url{https://en.wiktionary.org/wiki/Category:xtm:Oaxaca,_Mexico}
\end{itemize}
\section{Religious and Community Resources}
\begin{itemize}
\item Joshua Project profile: \url{https://joshuaproject.net/people_groups/14704}
\end{itemize}
\section{Digital and Audio Resources}
\begin{itemize}
\item Global Recordings Network, Mixtec, Magdalena Penasco (includes a Bible League New Testament from 2013 and a Jesus Film version, per the search summary): \url{https://globalrecordings.net/en/language/xtm}
\item Global Recordings, San Mateo variety: \url{https://globalrecordings.net/en/language/6443}
\item Global Recordings, Tlacotepec variety: \url{https://www.globalrecordings.net/en/language/4068}
\item Global Recordings, San Crist\'obal Amoltepec variety: \url{https://www.globalrecordings.net/en/language/25419}
\end{itemize}
\section{Curriculum Development Context}
With several thousand speakers there is a real base, but the search found no teaching materials. Dialect variation across the four towns would need to be addressed in any orthography and curriculum.
\end{document}